\section{Private sign bridges}
\label{sec:private-bridges}

A private sign bridge is an $N$-person group that can co-sign files. Each
node is its own storage: \code{M.A.1} is an employee, \code{M.A} is the
accounting manager, \code{M} is a cross-department manager (a CXO). A
bridge of \code{M.S.2}, \code{M.S.3}, and \code{M.A.2} therefore notifies
\textbf{five} stores --- those three employees plus department managers
\code{M.S} and \code{M.A}. The CXO is not in that set unless a member is
itself a department manager.

Members receive a sealed copy of a shared Ed25519 key (plus salts).
Managers receive roster metadata only, so they can track the live standard.
This database never keeps another person's sealed secret; \code{create} and
\code{remove-member} write one \file{.kqpb} \textbf{envelope} per store.

\begin{noteBox}
\code{create} and \code{remove-member} use the same plan-then-commit
pattern: they generate every envelope first, write the files, and only then
persist the new (or rotated) bridge in this database. Either both land or
neither does --- a failed write leaves no live row for that change, and a
failed commit deletes the envelopes it just wrote, since they would
describe a bridge this store never recorded, and \file{.kqpb} files are
never overwritten (leftovers would block the retry). Either way, simply run
the command again. There is no separate ``redeliver initial packages''
command because a failed create never commits.
\end{noteBox}

Each \flag{--member} must already have a registered \textbf{signing} public
key under that label (\code{generate --type signing --label M.S.2
--register}). The roster stores that key. \code{verify} checks the artifact
against the roster-bound personal key, not a public key declared only on
the signature.

\subsection{Envelope shape}

Each \file{.kqpb} is a cryptographic envelope (the same idea as a
\code{KQXB} export bundle): the outside names the recipient's X25519 public
key; the inside is \code{crypto_box}-sealed so only that store's
encryption private key can open it. A mailbox, USB drop, or email can carry
the envelope without being able to read the letter. (\file{.kqbn} eviction
notices are the exception --- they are public routing slips, not sealed
envelopes.)

\begin{center}
\begin{tabular}{@{}p{0.46\linewidth}p{0.46\linewidth}@{}}
\toprule
\textbf{outside (anyone can see)} & \textbf{inside (recipient only)} \\
\midrule
KQPB magic, kind & wrap\_salt \verb!||! bridge secret (members) \\
recipient encryption pub & or roster + salts (managers) \\
 & roster includes each member's personal signing pub \\
 & + rotate/destroy auth signature \\
\bottomrule
\end{tabular}
\end{center}

Five people in the example above means five envelopes, each addressed to a
different public key. Operators can copy those files out of band, or push
them through the mailbox relay (\code{keyquorum relay push}) so each store
can \code{relay pull --import} locally.

\begin{lstlisting}
# Each member needs a registered signing public key under their label:
keyquorum generate --type signing --label M.S.2 --register \
  --public-key-out M.S.2.sign.pub > M.S.2.sign.key
# ...same for M.S.3 and M.A.2

# On any machine that can see the tree (and manager encryption pubs):
keyquorum bridge private create 1 \
  --member M.S.2 --member M.S.3 --member M.A.2 \
  --supervisor M.S=SoftwareManager.pub --supervisor M.A=AccountingManager.pub \
  --self M.S.2 --output-dir ./bridge-packages --label eng-acct

# Each other person imports into *their* database:
keyquorum --db ma2.sqlite bridge private import \
  --file ./bridge-packages/M.A.2.kqpb --share-file Accounting2.key

# Sign with the bridge key + the member's personal signing key:
keyquorum sign --bridge-uid <uid> --node M.A.2 \
  --signing-key-file Accounting2.sign.key --share-file Accounting2.key \
  --message-file report.pdf --signature-out report.kqbs

# Peer verifies with their membership and the artifact:
keyquorum --db ms2.sqlite verify --bridge-uid <uid> --as-node M.S.2 \
  --message-file report.pdf --signature-file report.kqbs
\end{lstlisting}

\subsection{Revocation and rotation}

If \code{M.S} revokes \code{M.S.3}'s hardware key, that label is dropped
from every live private bridge. Ordinary \code{revoke} (not only
\flag{--evict}) prints the notify list and writes a \file{.kqbn} notice
\textbf{after} the hardware revoke commits. Remaining members must rotate
--- the departed person still holds the old bridge secret. A remaining
member then runs:

\begin{lstlisting}
keyquorum bridge private remove-member <uid> --member M.S.3 \
  --node M.S.2 --share-file Software2.key --output-dir ./bridge-packages
# Copy the new packages to M.S.2, M.A.2, M.S, M.A, and M.S.3
# (or `keyquorum relay push --dir ./bridge-packages`).
\end{lstlisting}

A two-person bridge is destroyed when one member is removed.

\subsection{Evicting a hardware key}

Banning a hardware key is \code{revoke <hardware-id>}. That always drops
pairings and whitelist rows for every live leaf sealed to that token, and
drops that label from every live private sign bridge. \flag{--evict}
additionally PSS-refreshes survivors of that leaf (\flag{--key-id} /
\flag{--node}, or the unique leaf the token backs). Eviction needs a parent
threshold of at least 2, and every remaining active sibling must be a
hardware-backed leaf --- a 1-of-$N$ parent is refused. Binds between
survivors stay. \flag{--share-file} is that sibling's actual key file
(\file{.pub} / \file{.key} from \code{generate} or \code{split
--generate-keys}, or a PEM / OpenSSH public key). A \file{.pub} file uses a
sibling \file{.key} when present.

\begin{lstlisting}
keyquorum split --label "team escrow" --threshold 2 \
  --leaf alice=alice.pub --leaf bob=bob.pub --leaf carol=carol.pub \
  --generate-keys --register
keyquorum tree <key-id>
keyquorum revoke <carol-hardware-id> --evict \
  --share-file alice.key --share-file bob.pub
# pin a leaf when the token backs more than one:
#   --key-id <key-id> --node carol
\end{lstlisting}
